\section{Limitations}
\label{sec:limitations}

\paragraph{Scope.} Typical makes text-only, bounded decisions over candidate sets that fit in a short
suffix: categorical, ordinal and Bernoulli contracts with abstention. It does not plan, search, call
tools or read images. Per-query cost grows with the rendered candidate tokens, so candidate sets in the
hundreds are served through a shortlisting front end (\Cref{sec:results:cost}).

\paragraph{Multi-step composition.} Decisions that require carrying an intermediate value through
several steps are the open frontier. On DecisionMix level 7 (temporal, numeric, expected-value and
trade-off composition, never trained) the v1 1.7B and 4B models score .493 and .544 against a .441 guess rate,
and the 14B model .620. On the JevBench hard tier a frozen 14B backbone read with three in-context
exemplars scores above our trained 14B checkpoint (paired difference $+.108$ $[+.027, +.189]$,
$p = .015$), while the trained model leads on the standard tier. \Cref{app:hard-tier} gives the
per-family breakdown. \Cref{sec:discussion} describes the extension this points to.

\paragraph{Statistics.} The readout comparison in \Cref{sec:analysis:readout} and the evidence-first
result in \Cref{sec:results:long} are replicated across seeds; other rows are single training runs.
JevBench numbers use the 231 public items and are not a ranked leaderboard entry. Its intervals are
$\pm6$--13 points on the standard tier (36 independent states) and $\pm9$ points on the hard tier
(111 items), so we base comparisons between our own checkpoints on paired tests and on larger
purpose-built sets.

\paragraph{Generalization and invariance.} ``Runtime-defined labels'' refers to the axes of
\Cref{tab:generalization}. The rubric-generalization results share one rule engine between training and
held-out families. Choice is order-robust rather than order-invariant: the contextual readout reduces
order sensitivity relative to the letter interface, per-step option shuffling during training
contributes to that robustness, and only Noul is exactly invariant. The pairwise-odds property of
\Cref{eq:factored-null} holds among non-null candidates; the abstention gate is set-dependent by design.

\paragraph{Ladder and release.} The size ladder is one recipe on one backbone family (Qwen3) with a
Qwen3.5-4B checkpoint as a second-generation point; the Qwen3-8B rung underperforms its neighbours in
both frozen and trained form. Inference code, weights and evaluation outputs are public; the training
code is not yet released.

\section{Conclusion}
\label{sec:conclusion}

Typical reads decisions directly out of a pretrained causal language model. A state is encoded once, a
question and its runtime candidates are read in a short suffix, and a head at mid-depth returns a typed
distribution with a factored abstention event, at 45\,ms per decision at 1.7B and 59\,ms at 14B. The design rests on measured answers to where the decision lives: inside the forward pass
over the candidates, in their contextual states, below the layers that specialize for generation. The
training recipe shows that decision types, counterfactual rubrics, abstention targets and selection
criteria shape a decision model as strongly as its architecture, and the render-crossed evaluation
shows how to measure such effects cleanly. The models and the inference package are public, and the
same interface is the natural terminal layer for decision systems that add intermediate computation
before the typed readout.
